\scanpage{35}
\mspage{15}

\uline{une} \uuline{définition} \uline{de la coh. $p$-adique indépendante de la coh. plate}, \uline{quitte} \uuline{ensuite} \uline{à étudier ses relations (sans doute} fort intéressantes) \uline{avec cette dernière}.

2) \uline{L'approche de Washnitzer-Monsky.}

C'est dans cette approche qu'on voit la cohomologie de De Rham \deleted{qui} montrer à nouveau le bout de l'oreille ! Nous allons expliquer \inserted{pour le moment,} comment elle s'introduit dans l'approche de W-M., en \unsure{simplifiant} leur construction à l'extrême, (au risque de la défigurer \ldots !). On part de $X_0$ \unsure{projective} lisse sur $s$, \uline{on suppose} \inserted{\uline{d'abord}} qu'on puisse relever en $X$ \deleted{projective} propre et lisse sur $S=\operatorname{Spec}W(k)$, \deleted{soit $\eta$ le point générique}. Considérons le complexe $\underline{\Omega}_{X/S}^{\bullet}$, et

\margin{\diagram{\begin{tikzcd}[ampersand replacement=\&]
X_0 \arrow[r,phantom,"\subset"] \arrow[d,dash] \& X \arrow[d,dash]\\
s:\operatorname{Spec}k \arrow[r,phantom,"\subset"] \& S:\operatorname{Spec}W
\end{tikzcd}}}

\[
\boxed{H^i_{\mathrm{DR}}(X/S)=\mathbb H^i(X,\Omega^\bullet_{X/S})}
\]
C'est un module de t.f. sur $W(k)$.

\margin{34}
